\section{Розділ~\ref{sec:dendrites}. Кількість дендритів}

Будову класів і кількість входів виміряно анатомічно й електричними записами; оцінка~\eqref{eq:dendriteload} --- проста фізика конденсатора, а обчислювальне пояснення кількості входів спирається на теорію й моделі.

{\footnotesize
\setlength{\tabcolsep}{3pt}\renewcommand{\arraystretch}{1.15}
\begin{longtable}{>{\raggedright}p{0.5cm}>{\raggedright}p{4.0cm}>{\raggedright}p{5.6cm}>{\raggedright}p{2.3cm}>{\raggedright\arraybackslash}p{1.7cm}}
\toprule
№ & Твердження & Дослід і що виміряно & Джерело & Статус \\
\midrule\endfirsthead
\toprule
№ & Твердження & Дослід і що виміряно & Джерело & Статус \\
\midrule\endhead
1 & Питома ємність мембрани $c_m\approx1$~мкФ/см$^2$ & Із тіла нейрона знімали мембрану на піпетку, подавали сходинку напруги й за спадом ємнісного струму знаходили повну ємність, потім ділили на площу: $0.9$~мкФ/см$^2$ у трьох класів нейронів & \cite{Gentet2000} & виміряно \\
2 & Час заряджання $t\approx c_mA\Delta V/I$~\eqref{eq:dendriteload}: великому нейронові потрібні десятки одночасних входів, маленькому вистачає одного великого синапсу & Оцінка за законом заряджання конденсатора; числа ($20$ і $300$~пФ, $0.1$ і кілька нА) --- порядки величин & виведення в розділі & теорія \\
3 & Один величезний синапс дає струм у кілька наноамперів і відразу доводить маленький нейрон до порога & Одночасний запис закінчення й отримувача у зрізі: струм одного синапсу $2$--$13$~нА, кожен імпульс входу викликає надпорогову відповідь, затримка близько $0.4$~мс за $36^\circ$C; отримувач видає \nt{GLYC} & \cite{Borst1995,BorstSoria2012} & виміряно \\
4 & Нуль дендритів: у чутливих нейронів дотику й болю імпульс іде від датчика в спинний мозок повз тіло клітини & Будову (аксон, що ділиться надвоє біля тіла) встановлено анатомічно. Те, що проведення не потребує збудливості тіла, показано на перевіреній моделі такого нейрона: без збудливого тіла імпульс проходить розгалуження так само надійно & \cite{AmirDevor2003} & теорія \\
5 & Один дендрит: нюховий нейрон несе один тип рецептора й передає одну вхідну лінію & Огляд дослідів із генами рецепторів: кожен нюховий нейрон експресує один ген рецептора з великої родини & \cite{Mombaerts2004} & виміряно \\
6 & Два дендрити: у детекторів напрямку на звук входи від лівого й правого вуха сідають на різні дендрити & Анатомія й записи в ссавців, зібрані в огляді: входи від двох вух рознесено по двох протилежних дендритах & \cite{Grothe2010} & виміряно \\
7 & Рознесення входів по двох дендритах робить «І» суворішим & Показано на моделі: насичення~\eqref{eq:shunt} на одному дендриті не дає входам з одного боку дійти до порога, і детекція збігів поліпшується. Прямого досліду, де змінювали б розкладку входів, немає & \cite{AgmonSnir1998} & теорія \\
8 & Близько $7\cdot10^{10}$ маленьких \nt{GLUT}-нейронів кола навчання рухів, по $\sim4$ входи в кожного & Підрахунок клітин ізотропним фракціонуванням: у цій частині мозку людини близько $69$ млрд нейронів. Запис однієї такої клітини в живого щура: на дотик надходять пачки дискретних струмів від небагатьох входів, і клітина спрацьовує, коли вони додаються & \cite{Azevedo2009,Chadderton2004} & виміряно \\
9 & Пороговий елемент з $n$ входами розділяє не більше $P_{\max}\approx2n$ випадкових образів~\eqref{eq:cover}; для вихідного нейрона кола навчання рухів це межа $\sim2\cdot10^5$, а реалістична оцінка --- $\sim4\cdot10^4$ відповідностей & Межа --- класичний результат комбінаторної геометрії. Реалістична оцінка --- теорія ємності лише з додатними вагами й запасом на шум, застосована до виміряної кількості входів цього нейрона & \cite{Cover1965,Brunel2004} & теорія \\
10 & Змішувачі з малою кількістю входів потрібні, щоб розширити вхід; «чотири» близько до оптимуму & Теорія: вимірність подання, яку дає шар змішувачів, максимальна приблизно за тієї кількості входів, що спостерігається анатомічно; вихідна гіпотеза розширення --- у класичних працях & \cite{Marr1969,Albus1971,LitwinKumar2017} & гіпотеза \\
11 & Великі дерева: $10^4$--$10^5$ входів & Підрахунок синапсів на електронних знімках: $\approx1.7\cdot10^5$ синапсів на одному вихідному \nt{GABA}-нейроні кола навчання рухів у щура; близько $30\,000$ збуджувальних і $1\,700$ гальмівних входів на \nt{GLUT}-нейроні гіпокампа & \cite{NapperHarvey1988,Megias2001} & виміряно \\
12 & Гілки великого дерева --- окремі підсхеми з власними порогами; нейрон --- маленька багатошарова мережа & Точкова стимуляція двох місць на тонких дендритах: входи на одній гілці додаються за сигмоїдою, на різних гілках --- лінійно. У дендритах нейронів кори людини знайдено градуальні кальцієві імпульси, що дають змогу одній клітині розв'язувати лінійно нероздільну задачу. Модель: щоб відтворити один нейрон, потрібна глибока мережа & \cite{Polsky2004,Gidon2020,Beniaguev2021} & виміряно \\
13 & Дендрити-виходи: кожна гілка \nt{GABA}-нейрона сітківки --- окремий детектор напрямку & Двофотонний запис кальцію в окремих гілках: сигнал у гілці вибірковий до руху від тіла клітини до кінчика, а напруга тіла --- ні & \cite{Euler2002} & виміряно \\
14 & Кількість входів іде за задачею (твердження~\ref{prop:dendrites}) & Будову класів виміряно (рядки 3--13); що кількість входів обрано заради швидкості чи класифікації, показано лише теорією й моделями для окремих класів & \cite{LitwinKumar2017,AgmonSnir1998} & гіпотеза \\
\bottomrule
\end{longtable}}
